\section{A complete adapted-noise argument}\label{sec:first-route}

Our first goal is to control a specified list of cards, uniformly over
its starting positions. We keep the deterministic coordinate schedule.
The conditional error is generated by transport across pairs with one
available position, and all earlier error is erased after one sweep.
We will first turn this finite memory into a scalar recurrence and then
substitute the resulting list estimate into a precise Fourier transfer.

\subsection{The exact noise and covariance identities}

Write $x^i=x+e_i$ for the coordinate-$i$ partner of $x\in V$, and
define $(P_i f)(x)=(f(x)+f(x^i))/2$. The operators $P_i$ commute,
and their product over all coordinates projects onto constant vectors.

Fix a collection of blocker labels and a different tagged label. The
starting labels may be deterministic, or may be sampled independently of
the switches and then revealed. Let $V_t$ be the positions unoccupied by
blockers, and write $m=|V_t|$. Given the revealed starting labels and the
blocker paths through time $t$, let $p_t$ be the conditional law of the
tagged position and define
\[
 f_t(x)=p_t(x)-\frac{\one_{V_t}(x)}m,
 \qquad a_t=\E\|f_t\|_2^2.
\]
Thus $f_t$ vanishes off $V_t$, sums to zero, and
$\|f_t\|_2^2=\|p_t\|_2^2-1/m\le1$.
This is the same type of centered second moment used in the
swap-or-not analysis~\cite[Section~3]{hmr2014}; here its evolution is
governed by the prescribed coordinate schedule.

By Lemma~\ref{lem:marginal-averaging}, the conditional vector is
computable from the blocker-path prefixes, even if the complete blocker
paths are given. Positions in $V_t$ are \emph{available}; all other
positions are occupied by blockers. On a pair with two available input
positions, the two entries are averaged. On a pair with one available
input position, its entry is transported to the unique available output
position. A pair with no available input position carries zero. The same
rule transports the uniform vector $\one_{V_t}/m$.

The same linear update can start from any deterministic real vector
$f_0$ supported on $V_0$ and satisfying $\sum_x f_0(x)=0$.
This more general vector need not be a difference of probabilities.
Every update preserves its sum and support, and averaging or transport
gives the pathwise bound
\begin{equation}\label{noise:pathwise-contraction}
 \|f_t\|_2^2\le\|f_0\|_2^2.
\end{equation}
The next two lemmas apply to this signed-vector evolution. For the
conditional tagged law above, $\|f_0\|_2^2=1-1/m\le1$.

\begin{lemma}[Finite memory of the centered noise]
\label{noise:memory}
For either the conditional tagged vector or the signed-vector evolution
just defined, put
\[
 \xi_t=f_{t+1}-P_{i_t}f_t,
 \qquad
 w_t=\E\sum_{x\in V}f_t(x)^2\one_{\{x^{i_t}\notin V_t\}}.
\]
Then $\xi_t$ is conditionally centered given the past, and
\begin{align}
 \E\|\xi_t\|_2^2&=\tfrac12w_t,\label{noise:noise-energy}\\
 a_t&=\sum_{j=1}^{d}2^{-j}w_{t-j}\quad(t\ge d),
 \label{noise:memory-equality}\\
 a_{t+1}&=\tfrac12(a_t+w_t-2^{-d}w_{t-d})
             \le\tfrac12(a_t+w_t)\quad(t\ge d).
 \label{noise:one-step-memory}
\end{align}
Moreover, if $C_t=\E[f_tf_t^{\mathsf T}]$ and $e_x$ is the unit vector
at $x$, then the exact covariance recursion is
\begin{align}
 C_{t+1}&=P_{i_t}C_tP_{i_t}\notag\\
 &\quad+\frac14\sum_{\{x,z\}:z=x^{i_t}}
  \E\!\left[\one_{\{|\{x,z\}\cap V_t|=1\}}
                     (f_t(x)+f_t(z))^2\right]
     (e_x-e_z)(e_x-e_z)^{\mathsf T},
 \label{noise:covariance}
\end{align}
where the last sum is over unordered pairs.
\end{lemma}

\begin{proof}
Use either the blocker-prefix filtration or the larger filtration
containing the initial selection and all switch coins already used.
The vector $f_t$ is measurable for either; the larger filtration does not
redefine the path-conditioned law $p_t$. On a pair with one available position $x$,
its contribution to $\xi_t$ is
$\pm f_t(x)(e_x-e_{x^{i_t}})/2$, with a fair sign. The signs for
different pairs are conditionally independent. On the other pairs the
noise is zero. This proves conditional centering,
\eqref{noise:noise-energy}, and \eqref{noise:covariance}; the cross term
between $P_{i_t}f_t$ and $\xi_t$ also has conditional mean zero.

Unroll $f_{s+1}=P_{i_s}f_s+\xi_s$ over the last $d$ updates. The initial
term is zero because its coordinate averages include every direction.
Noises created at different times are orthogonal in expectation, even
after the subsequent deterministic operators, by conditional centering.
A pair difference created at time $t-j$ is followed by $j-1$ averages
in distinct other directions. Its endpoints spread over two disjoint
subcubes of size $2^{j-1}$, so its squared norm is multiplied by
$2^{-(j-1)}$. Different pair noises at one time have zero cross terms
in expectation by their independent signs, even if the propagated
supports overlap. Consequently
\[
 a_t=\sum_{j=1}^d2^{-(j-1)}\E\|\xi_{t-j}\|_2^2,
\]
which proves \eqref{noise:memory-equality}. Subtracting one half of that
identity from its version at $t+1$ gives
\eqref{noise:one-step-memory}. Equivalently, taking traces after
expanding \eqref{noise:covariance} gives these same weights: the factor
$1/4$ multiplies the squared pair-difference norm
$2\cdot2^{-(j-1)}$, giving exactly $2^{-j}$.
\end{proof}

The next elementary comparison will be used repeatedly. If for all
$s\ge s_0$ one has $w_s\le\rho a_s+\eta$, with $0\le\rho<1$, then
\eqref{noise:memory-equality} implies
\begin{equation}
 a_t\le\rho\sum_{j=1}^d2^{-j}a_{t-j}+\eta
 \qquad(t\ge s_0+d).
 \label{noise:comparison-recurrence}
\end{equation}
For any $1/2<\gamma<1$ with $\rho/(2\gamma-1)\le1$, the comparison
\begin{equation}
 a_t\le a_0\gamma^{t-t_0}+\frac{\eta}{1-\rho},
 \qquad t_0\ge s_0+d,
 \label{noise:comparison-solution}
\end{equation}
holds for all $t\ge0$: it is immediate up to $t_0$ from $a_t\le a_0$;
after that, induction uses
$\rho\sum_{j\ge1}(2\gamma)^{-j}=\rho/(2\gamma-1)$ and
$\sum_{j=1}^d2^{-j}\le1$.

\subsection{Uniform marginals from the preceding coordinate split}

We first keep the initial blockers completely arbitrary. The previous
use of a coordinate makes the blocker counts in its two halves nearly
equal, and the intervening layers separate the two sources of randomness.

\begin{lemma}[Half-density bound for arbitrary blockers]
\label{noise:halves}
Fix $0<\alpha\le1$, assume that at least $\alpha n$ positions are
available, and let $f_0$ be any deterministic sum-zero vector supported
on those positions. Set
\[
 \delta=\frac\alpha2,\qquad
 \eta_n=2\exp\!\left(-\frac{\alpha^2n}{4}\right),\qquad
 \gamma=1-\frac\alpha4.
\]
Uniformly over deterministic starting positions and choices of labels,
\begin{equation}
 \begin{split}
 w_s&\le(1-\delta)a_s+\eta_n\|f_0\|_2^2\quad(s\ge d),\\
 a_t&\le\left(\gamma^{t-2d}+\frac{\eta_n}{\delta}\right)
                  \|f_0\|_2^2\quad(t\ge0).
 \end{split}
 \label{noise:half-energy}
\end{equation}
\end{lemma}

\begin{proof}
The preceding update in direction $i_s$ is $s-d\to s-d+1$.
Conditional on the blocker configuration immediately before that update,
the blocker count in one coordinate half afterwards is a constant from
two-blocker pairs plus a sum of at most $n/2$ independent fair Bernoulli
variables from mixed pairs. Its mean is half the total blocker count.
If either half has blocker fraction greater than $1-\alpha/2$, its count
differs from this mean by at least $\alpha n/4$. Hoeffding's bound therefore
gives probability at most $2e^{-\alpha^2n/4}$ for this event.

Condition now on the blocker paths just after that preceding update.
The next $d-1$ updates use all other coordinates, acting independently
in the two halves. The value $f_s(x)$ is computed from the history in
the half of $x$ only. Its opposite partner $x^{i_s}$ is in the other
half. Every blocker in that other half has uniform one-card marginal
there after these $d-1$ coordinate updates, because each bit is refreshed
once. Thus the partner-blocker indicator is conditionally independent
of $f_s(x)^2$, and its conditional mean is that half's earlier blocker
fraction. On histories with both fractions at most $1-\delta$, summing
over $x$ bounds the conditional contribution by $1-\delta$ times the
conditional energy. On the remaining histories use
$\|f_s\|_2^2\le\|f_0\|_2^2$.
This proves the first inequality. Apply
\eqref{noise:comparison-solution}, with $\rho=1-\delta$,
$t_0=2d$, and $2\gamma-1=1-\delta$, to obtain the second.
\end{proof}

\begin{proposition}[Two uniform large-list estimates]
\label{noise:uniform-marginals}
For every sufficiently large $d$, the following estimates hold uniformly
over specified ordered lists and deterministic starting decks.
\begin{enumerate}
\item At time $1024d$, the images of any at most $15n/16$ labels have
total-variation distance at most $n^{-3/2}$ from a uniform injective
tuple.
\item At time $256d$, the images of any at most $7n/8$ labels have
total-variation distance $o(1)$ from a uniform injective tuple.
\end{enumerate}
\end{proposition}

\begin{proof}
For the first assertion use $\alpha=1/16$ in Lemma~\ref{noise:halves}.
At $T=1024d$,
$\gamma^{T-2d}\le e^{-1022d/64}$, so $a_T\le2n^{-6}$ eventually.
For the second use $\alpha=1/8$, giving $\gamma=31/32$ and
\[
 a_{256d}\le e^{-254d/32}+16\eta_n,
 \qquad n^3a_{256d}=o(1).
\]
In each case expose the list in order. For its next card take the
preceding labels as blockers. Conditional on their complete paths, its
distance from uniform on the remaining positions is
$\|f_T\|_1/2$, whose expectation is at most $\sqrt{na_T}/2$.
Conditioning on just their endpoints can only decrease this expected
distance, by convexity, since the uniform reference on available positions depends
only on those endpoints. Interpolate between laws with a true prefix
and a suffix extended uniformly without replacement. Adjacent laws
differ by at most the preceding expected one-card error; summing at
most $n$ errors gives $n^{3/2}\sqrt{a_T}/2$. The first bound is at most
$n^{-3/2}$ and the second tends to zero. This proves both assertions.
\end{proof}

\paragraph{Uniform entropy as a consequence.}
The same estimate gives more than total variation. Fix
$0<\alpha<1$ and $A>0$. There is an integer $C=C(\alpha,A)$ such
that, for every specified ordered list $(a_1,\ldots,a_k)$ with
$k\le(1-\alpha)n$, and uniformly over deterministic starting decks,
\begin{equation}\label{noise:uniform-entropy}
 D\!\left(\mathcal L((\Pi_{Cd}(a_j))_{j=1}^k)
       \,\middle\Vert\,\Unif(\operatorname{inj}([k],V))\right)
 =O(n^{-A}).
\end{equation}
Indeed, for a probability vector $p$ on $m$ points,
$\log u\le u-1$ gives
$D(p\Vert\Unif_m)\le m\|p-\Unif_m\|_2^2$.
The entropy chain rule sums this bound for the successive conditional
cards. Conditioning further on the preceding paths can only increase
the mean entropy, by convexity, since the uniform reference depends
only on their endpoints. With $m_j=n-j+1$ and $f_T^{(j)}$ the
path-conditioned error for the $j$th card, the tuple entropy is at most
\[
 \sum_{j=1}^k m_j\E\|f_T^{(j)}\|_2^2
 \le n^2\left(\gamma^{T-2d}+\frac{\eta_n}{\delta}\right),
 \qquad \gamma=1-\frac\alpha4,\quad\delta=\frac\alpha2,
\]
by Lemma~\ref{noise:halves}. Choose $C$ so that
$\gamma^{C-2}<2^{-(A+2)}$ and put $T=Cd$; the exponential
$\eta_n$ term is negligible. Thus later entropy calculations provide
alternative contraction mechanisms and more detailed intermediate
identities; small entropy for a fixed list is already a consequence
of the half-density argument.

\subsection{Conventions and the first full-deck application}
\label{sec:probability-completion}

The scalar calculation has now proved the two list assertions of
Theorem~\ref{thm:first-main}. The remaining step recovers the joint
order of the omitted labels as well. We first state the probability
completion used here and in later applications, and then give the exact
isotypic input for this route. Completions involving subprobabilities,
opposite orientations, kernels or different factor laws retain their
separate arguments at the point of use.

We use the characteristic-zero irreducible representations of $S_n$,
indexed by partitions $\lambda\vdash n$ and realized unitarily.
Write $\lambda_1$ for the first-row length and $\lambda'_1$ for the
first-column length. The row $(n)$ gives the trivial representation,
and the column $(1^n)$ gives sign; for $n\ge2$ these are exactly the
dimension-one representations. The standard background is in
\cite{sagan,etingof}.

\paragraph{Probability completion.}
For a measure $h$ on $S_n$ and an irreducible unitary representation
$\rho_\lambda$ of dimension $D_\lambda$, use the mass transform
$\widehat h(\lambda)=\sum_g h(g)\rho_\lambda(g)$ and the ordinary,
unnormalized Hilbert--Schmidt norm. Fourier matrices multiply in
convolution order. For every probability law $\sigma$, finite-group
Plancherel gives
\begin{equation}\label{eq:c-plancherel}
 4\|\sigma-U_{S_n}\|_{\TV}^2
 \le\sum_{\lambda\ne(n)}D_\lambda
              \|\widehat\sigma(\lambda)\|_{\HS}^2.
\end{equation}
This is the finite-group upper-bound method of Diaconis and
Shahshahani~\cite[Section~2 and Lemma~14]{diaconis-shahshahani1981}.
Indeed Cauchy--Schwarz bounds the left side by
$n!\sum_g|\sigma(g)-1/n!|^2$, which Plancherel identifies with the
right side.

Suppose $\mu,\nu$ are probability laws with
$\|\mu-\nu\|_{\TV}\le\delta_n=o(1)$ and
$|\widehat\nu(\mathrm{sgn})|=o(1)$. Let $p\ge1$ be a fixed integer,
independent of $n$. Setting $\sigma=\nu^{*p}$ gives
\begin{equation}\label{sweeps:completion}
 4\|\nu^{*p}-U_{S_n}\|_{\TV}^2
 \le\sum_{\lambda\ne(n)}D_\lambda
              \|\widehat\nu(\lambda)^p\|_{\HS}^2.
\end{equation}
An operator bound
$\|\widehat\nu(\lambda)\|_{\op}\le K D_\lambda^{-\alpha}$
gives the nonlinear summand bound
\[
 D_\lambda\|\widehat\nu(\lambda)^p\|_{\HS}^2
 \le K^{2p}D_\lambda^{2-2p\alpha}.
\]
A squared Hilbert--Schmidt bound
$\|\widehat\nu(\lambda)\|_{\HS}^2\le K D_\lambda^{-\beta}$
instead gives
\[
 D_\lambda\|\widehat\nu(\lambda)^p\|_{\HS}^2
 \le K^pD_\lambda^{1-p\beta}.
\]
Here $K,\alpha,\beta$ are fixed for the application. An $n$-dependent
prefactor that is uniformly bounded for all sufficiently large $n$ may
be replaced by a fixed upper bound $K$. The first estimate
uses $\|M\|_{\HS}\le\sqrt{D_\lambda}\|M\|_{\op}$; the second
uses Hilbert--Schmidt submultiplicativity. Neither requires normality.
If the resulting sum over $D_\lambda>1$ tends to zero, the sign
assumption and \eqref{eq:c-plancherel} give
$\|\nu^{*p}-U_{S_n}\|_{\TV}=o(1)$. Replacing $p$ factors one at a
time costs at most $p\delta_n$, by contraction under convolution with
a probability law. Hence $\mu^{*p}$ has the same limit.

In our applications $\mu$ is the position-increment law of a block of
$T$ physical shuffles, where $T$ is a positive integer divisible by
$d$. Independent such blocks contain whole sweeps and have product law
$\mu^{*p}=q_d^{*(pT)}$, by \eqref{eq:frames:cyclic}. Each block has
zero expected sign: changing one fair switch while fixing every other
coin changes its endpoint permutation by a transposition. Thus a nearby
probability $\nu$ has
$|\widehat\nu(\mathrm{sgn})|\le2\delta_n$.
If $\nu$ is obtained by conditioning on a retained event of probability
$z>0$, one also has the bound $(1-z)/z$, since the unconditioned sign
expectation is zero. These statements hold also for an event on switch
settings rather than on endpoints. Equation~\eqref{eq:worst-state}
then supplies uniformity over deterministic starting decks.

\paragraph{The fixed-list application.}
We use the following exact case of the isotypic transfer proved in
\cite[Lemma~\ref*{l-l:trim} and Proposition~\ref*{l-l:isotypic}]{partialFourier}.
We state its data explicitly to distinguish it from the operator-norm
transfers used later.

Partition the labels into $b$ disjoint nonempty blocks and let $H_i$
permute block $i$ while fixing its complement. For a nonnegative
measure $h$ of mass at most one on $S_n$, suppose
\[
 h(gH_i)\le B/[S_n:H_i]\qquad(g\in S_n,\ 1\le i\le b).
\]
With the Fourier convention just fixed, the transfer gives
\begin{equation}\label{eq:first-isotypic}
 \|\widehat h(\lambda)\|_{\HS}^2
 \le bBC_2D_\lambda^{-1+4/b},\qquad
 C_2=\sup_{m\ge1}\sum_{\tau\vdash m}D_\tau^{-2}<\infty.
\end{equation}
All restriction multiplicities are included, and the trace defining
$\|\cdot\|_{\HS}$ is not normalized. The same companion proves
\begin{equation}\label{eq:first-degree-sum}
 \sum_{\lambda\vdash n:\,D_\lambda>1}D_\lambda^{-2}\longrightarrow0
\end{equation}
in \cite[Lemma~\ref*{l-l:degree-reciprocal-two}]{partialFourier}.

Apply the $256d$ assertion to the complements of eight equal blocks,
and let $\mu=q_d^{*(256d)}$. The sum $\delta_n$ of their eight
marginal errors tends to zero. The endpoint tuple on a complement
specifies precisely a left coset $gH_i$: two permutations agree on
that complement exactly when they differ by multiplication on the
right by $H_i$. Delete every coset whose mass exceeds twice its
uniform mass, simultaneously for all eight subgroups. A deleted coset
$C$ satisfies $\mu(C)\le2(\mu(C)-U(C))$, so the total deleted
mass is at most $2\delta_n$. Let $z$ be the retained mass and
$\nu$ the normalized retained law. Then
\[
 z\ge1-2\delta_n,\qquad
 \|\nu-\mu\|_{\TV}=1-z=o(1),\qquad
 \nu(gH_i)\le\frac{2}{z[S_n:H_i]}\le\frac4{[S_n:H_i]}
\]
for sufficiently large $n$. The same retained law satisfies every cap.
Equation~\eqref{eq:first-isotypic}, with $b=8$ and $B=4$, therefore
gives $\|\widehat\nu(\lambda)\|_{\HS}^2\le32C_2D_\lambda^{-1/2}$.

Use the squared Hilbert--Schmidt case of the probability completion
with $K=32C_2$, $\beta=1/2$ and $p=8$. The contribution of
$D_\lambda>1$ for $\sigma=\nu^{*8}$ is at most
\[
 \sum_{D_\lambda>1}D_\lambda
       \|\widehat\nu(\lambda)^8\|_{\HS}^2
 \le(32C_2)^8\sum_{D_\lambda>1}D_\lambda^{-3}=o(1).
\]
The final equality follows from~\eqref{eq:first-degree-sum}.

The original block law has zero expected sign, so
$|\widehat\nu(\mathrm{sgn})|\le2\|\nu-\mu\|_{\TV}=o(1)$.
The probability completion therefore proves
$\|\nu^{*8}-U_{S_n}\|_{\TV}\to0$ and bounds the replacement cost by
$8\|\nu-\mu\|_{\TV}=o(1)$.
The original eight blocks take $8\cdot256d=2048d$ physical shuffles,
and~\eqref{eq:worst-state} makes the estimate uniform over all
deterministic starts. This completes Theorem~\ref{thm:first-main}.
